\section*{Capítulo II, Problema VIII}

\textbf{Problema:}

8. Encuentra la ecuación del plano que pasa por los siguientes tres puntos.

\begin{itemize}
    \item[(a)] $(2,1,1)$, $(3,-1,1)$, $(4,1,-1)$
    \item[(b)] $(-2,3,-1)$, $(2,2,3)$, $(-4,-1,1)$
    \item[(c)] $(-5,-1,2)$, $(1,2,-1)$, $(3,-1,2)$
\end{itemize}

\textbf{Solución:}

Si $A$, $B$ y $C$ son tres puntos del plano, entonces dos vectores contenidos en él son $\overrightarrow{AB}$ y $\overrightarrow{AC}$. Un vector normal se obtiene con el producto cruz $\overrightarrow{AB} \times \overrightarrow{AC}$.

\begin{itemize}
    \item[(a)] Tomemos
    \[
        A=(2,1,1), \quad B=(3,-1,1), \quad C=(4,1,-1).
    \]
    Entonces
    \[
        \overrightarrow{AB} = (1,-2,0), \qquad \overrightarrow{AC} = (2,0,-2).
    \]
    Un vector normal es
    \[
        \overrightarrow{AB} \times \overrightarrow{AC} = (4,2,4) = 2(2,1,2).
    \]
    Podemos usar $N=(2,1,2)$. La ecuación del plano es
    \[
        2(x-2) + (y-1) + 2(z-1) = 0,
    \]
    es decir,
    \[
        2x + y + 2z - 7 = 0.
    \]

    \item[(b)] Tomemos
    \[
        A=(-2,3,-1), \quad B=(2,2,3), \quad C=(-4,-1,1).
    \]
    Entonces
    \[
        \overrightarrow{AB} = (4,-1,4), \qquad \overrightarrow{AC} = (-2,-4,2).
    \]
    Un vector normal es
    \[
        \overrightarrow{AB} \times \overrightarrow{AC} = (14,-16,-18) = 2(7,-8,-9).
    \]
    Usando $N=(7,-8,-9)$, obtenemos
    \[
        7(x+2) - 8(y-3) - 9(z+1) = 0,
    \]
    o equivalentemente,
    \[
        7x - 8y - 9z + 29 = 0.
    \]

    \item[(c)] Tomemos
    \[
        A=(-5,-1,2), \quad B=(1,2,-1), \quad C=(3,-1,2).
    \]
    Entonces
    \[
        \overrightarrow{AB} = (6,3,-3), \qquad \overrightarrow{AC} = (8,0,0).
    \]
    Un vector normal es
    \[
        \overrightarrow{AB} \times \overrightarrow{AC} = (0,-24,-24).
    \]
    Podemos tomar $N=(0,1,1)$. Así, la ecuación del plano es
    \[
        (y+1) + (z-2) = 0,
    \]
    esto es,
    \[
        y + z - 1 = 0.
    \]
\end{itemize}